\documentclass[11pt]{article}
\usepackage[margin=0.9in]{geometry}
\usepackage{amsmath,amssymb,hyperref}
\newcommand{\nm}[1]{\nolinkurl{#1}}
\emergencystretch=2em
\newcommand{\al}{\mathrel{\sim}}
\setlength{\parskip}{.45em}\setlength{\parindent}{0pt}
\title{Independent review of the ambient four-to-two exchange}
\author{Contextual Composition, issue 105}\date{20 November 2026}
\begin{document}\maketitle
\paragraph{Audit verdict.}
Accept the bounded claim: \nm{AmbientFourToTwo.exchange} proves the parent's exact pre-final exchange demand, without assuming its terminal tests. Independent native compilation and a 3,750-name axiom audit cover the frozen source. In one concrete thick-payload example the two actual packed tips are noninvertible, the exchange is invertible, and flipping its terminal arrow fails despite preserving its entire grade observation. The one material repair is to the replay instructions: a normal checkout must be filtered by the 259-file parent manifest before invoking the author's exact-membership assembler. Neither horizontal coherence as a whole nor representable recovery is certified.
\section*{Interface and proof audit}
The claim is pointwise in every terminal presheaf decoration, not equality of grade objects or of composites surviving one more multiplication. On strict category-valued presheaves over an ordinary base with split grade system, take the established ambient bicategory of strong transformations and modifications. Fix a native multiplication $m$ (including preservation of cartesian grade arrows), not a native package of triangle/pentagon witnesses. Set $M=M_X$. In the author's notation the ambient modification is exactly $(\mathtt{squareM}\ S\ m\ M)^{-1}$: its flat $4\to2$ source reads $\operatorname{wide}(m);\operatorname{inside}_1(m)$, its flat target reads $\operatorname{inside}_2(m);\operatorname{suffix}(m)$. \,No external assumption supplies the remaining terminal equation. \nm{Law} is an abbreviation inherited byte-for-byte from the parent. In it, $v\al w$ means that equality transports at both endpoints align the arrows, not conjugacy by arbitrary isomorphisms.
The row test has three layers. Its head-grade map is a transport. The observed last-grade map must align with a grade arrow $w$ having both its prescribed base and the equation $w;r=l$ with the \emph{original} right and left cones. Independently the $X$-payload is a transport. The law identifies the whole decorated component, before the functor that would postcompose it with multiplication; the whiskered conclusion uses that functor only in the forward direction.
\paragraph{Terminal calculation, not checklist laundering.}
For general routes $f,g:2\to1$ with cartesian-preserving grade part of $g$, write $u$ for the decorated wide chosen restriction, $r$ for the snoc chosen restriction in the inside frame, $l=\operatorname{raw}(g)(u)$, and $c$ for the final decorated leaf of the native comparison. Read \nm{AmbientTerminalCones.rawFactor} first. It aligns the grades of $c;r$ and $l$ by \nm{assembledTip} and \nm{exchangeFac}, referring to \nm{rightCone}/\nm{leftCone}. The two payloads are pure: on $l$ by functorial restriction plus its naturality transport; on $c;r$ by native payload purity plus the right restriction's chosen payload. Equality of target objects is necessary and comes from \nm{innerTail}, a full compositional frame. Chosen-arrow recognition therefore yields $c;r\al l$. This part alone says nothing about the ambient square.
\paragraph{Why these are the weak square's actual tips.}
Put $D_n$ for the chart taking a raw decorated length-$n$ row to its packed description. \nm{inputTip} identifies $D_2(u)$ with the weak square's input tip. The base is $f$'s leg on the original, $2+2$-split prefix (\nm{moveBase}, \nm{inIndex}); the target is \nm{inputLeaf}. Its packed fibre component is pure, not merely invertible: \nm{packedTwo} uses the $1+1$ stage, its head and decorated leaf frames, and the one-layer chosen-arrow calculation twice. \nm{outputTip} separately identifies $D_1(r)$ up to the serial source transport with the output tip. Here the target comes from \nm{outputLeaf} and the inside frame under the operation $g$, the base from \nm{snocBase} transported to that frame, and the packed purity from \nm{packedSnoc}. Neither identification assumes a strong map sends an arbitrary chosen lift strictly to a chosen lift.
\paragraph{Cancellation and return.}
For $a=\operatorname{oper}(g)$ the nested ambient mate $s$ satisfies $s;t_{\rm out}\al\operatorname{tot}(a)(t_{\rm in})$ by \nm{AmbientMateCalculus.factor}, an existing theorem about the selected weak comparison. The first tip and the $g$-operation chart identify the right side with $D_1(l)$. The second tip and the native factorization identify this with $D_1(c);D_1(r)$. \nm{AmbientNativeTerminal.finish} cancels only the original cartesian $t_{\rm out}$ and only over equal projected maps; purity of the two projections plus source alignment supplies that equality. There is no assumption of a monic or invertible directed leg. In \nm{readDiagram}, changing from a flat input to its $2+2$ split and staging the output $1+1$ compares evaluations of the \emph{same} total modification. The head-grade transport on each side comes from its own grade law, while the decorated terminal equality is the cancellation just established. The fixed full endpoint frames permit reassembly by the two pullback projections; the packed chart has a strict inverse and reflects aligned rows. Specialize $f=g=m$: \nm{muRead} supplies the genuine nested reading at $M$, not a hypothesized flat reading. Only thereafter do \nm{flatPoints}, \nm{exchange} and \nm{row} invoke the old checklist equivalence. This detour is logically redundant after \nm{current} but not circular: the proof of \nm{current} has no dependency on \nm{flatPoints} or \nm{exchange}.
\section*{Horizontal boundary}
The four-to-two proof uses the actual multiplication at $M_X$. By contrast, \nm{AmbientRightWeight.read} quantifies over every ambient strong arrow $f$, but its object is specifically the inverse of the folded inside-unit boundary $K(\eta_f)$, read in arities $1\to2$ at the native right-unit letter. \nm{AmbientUnitSquare.naturality} then accepts a $1\to2$ operation, a cell from its composite with $m$ to the identity, the two operation frames and a square that reads the operation at $f$. Its result is a square for the \emph{reframed} native cells. These neither assume nor prove arbitrary remaining generator readings. The exact public \nm{AmbientHorizontalNaturality.Right} specialization does not occur among the candidate declarations. Neither the public left/associativity squares, the complete horizontal triangle/pentagon, a packaged ambient pseudomonad, nor Kleisli recovery follows from the four-to-two reading alone. In particular the terminal $X$-transport is not whole-crossing purity; nor does it say the leg through which cancellation ran was invertible.
\section*{Source crosswalk and transport limit}
The relative-construction manuscript's \nm{eq:groth} defines the decorated fibre, with a grade map and a morphism to a reindexed terminal object. Its \nm{eq:grothmonad} defines multiplication by $P\otimes R$ together with restriction of the decoration along $\delta_{P,R}$. It explicitly does not invert $\delta$. The four-to-two result certifies a component needed when comparing two disjoint uses of this multiplication on actual arrows; it is a derived coherence obligation, not a separately labelled exchange theorem in the manuscript. \nm{prop:groth}, in contrast, asks for a composition-compatible local equivalence on representables. No such equivalence, nor an equivalence on the entire presheaf base, is part of this certificate. The exact source excerpt at the commit named by issue 105 is in the evidence archive.
\section*{Concrete discriminators}
The moving test and the payload test can be run \emph{on one actual row}. Use the natural-number preorder as base, the two distinct shades with a unique isomorphism between them, and the one-object terminal groupoid with Boolean arrows under exclusive-or. At the four-grade input of \nm{MovingAmbientExchange}, the left last-grade shade is white, the right black. Projecting a hypothetical equality of native endpoints to grades contradicts this, and the two original functor frames carry the inequality to the ambient endpoints. \nm{MovingThick} now specializes the proved law to this nonthin presheaf and this input. The actual arrow and its postcomposition with the nonidentity Boolean flip are both invertible and have the same endpoints and grade map: both pass all the grade tests. The actual bit is false (identity); the altered bit is true, so the latter arrow fails the $X$-payload test. They are unequal without comparing objects as a surrogate for arrows.
\paragraph{Directed witnesses in that same row.}
The original prefix leg is $0\to4$; an inverse would give $4\le0$. The projected bases of the actual packed input and output tips align with this same leg, via the split prefix and preservation of grade by the endpoint action. A functor preserves invertibility, including after equality transport, so neither actual tip is invertible. Both \emph{native} fixed cones are likewise noninvertible: project the left cone using the native move's base formula, and use \nm{coneBase} for the right. Thus the cartesian output tip that the proof cancels really can be noninvertible, not merely unproved invertible. See \nm{bothActualTips} and \nm{bothNativeCones} (namespace \nm{ExchangeReview.Directed}).
\paragraph{Wrong grade with pure payload.}
Independently, in the signed ordered-grade fixture the counterfeit in the fixed ambient frames is invertible and has pure terminal payload but the wrong terminal cone factor. \nm{StrongGrade} turns the author's three separate observations into a single formal negation: the actual ambient component does \emph{not} align with the conjugated counterfeit. Indeed, alignment of the two, composed with the counterfeit-to-native and actual-to-native alignments in the appropriate directions, would align the known wrong native arrow with the required one. The passing grade-only counterfeit in the Boolean model shows why this test cannot replace the independent terminal-payload check.
\paragraph{Cancellation pressure.}
In the official Boolean-endofunction test the same nonmonic constant map under two projections separates the cancellation hypotheses. Under terminal projection the base maps agree and composites agree but the leg is not cartesian. Under identity projection the leg is cartesian and composites agree but bases disagree. The competitors in both tests are unequal invertible arrows. Cartesian cancellation over a \emph{fixed equal base} survives both controls; general right cancellation does not.
\section*{Kernel, reproducibility and disposition}
All 259 canonical parent files independently match the native Git blob IDs at merge 230 (\nm{22f967a29c13b5bddea550a3b3ccf4f24b94dfae}). No missing historical check was supplied from a neighbouring project. The 295-file target is the parent plus 35 production modules and one official control, with only the umbrella, ordered audit and theorem index changed among existing sources. All 270 production modules are reachable from the umbrella. None imports an independent control. The author ZIP was checked against its 207-entry member digest list (208 members including the list).
A fresh project with no project build products used Lean 4.35.0 (\nm{eb6cafb7290e8bd5f2937522d8d804f58f02bafc}) and the nine pinned external dependency checkouts, including mathlib \nm{56a005855ea17e0b0958f73ec801ae140b0a1679}. The clean build took 42 minutes 45 seconds (04:15:24--04:58:09 UTC, 20 November 2026); 271 actual project compile records include the umbrella, whereas Lake's 1,018 jobs include dependencies and non-project tasks. For each compiled module, the source path, nonsynthetic compiler trace, native setup, and five product hashes were checked. The 89,340 project import artifact references (847,176 references overall; repeated per consumer) all resolved to this project's fresh outputs; other imported artifacts resolved to the pinned upstream or native toolchain roots.
The independent axiom transcript matches the sealed author's transcript byte for byte. Among 3,750 distinct names printed in index order, \nm{propext}, \nm{Classical.choice}, and \nm{Quot.sound} occur in respectively 3,531, 3,364 and 3,424 dependency lists; among the 129 additions the counts are 128, 126 and 126. No other axiom occurs. All 17 official compiler probes passed. A process interruption during the seventh probe required a verified-source restart of that probe and the remaining ten, without restarting or reusing a different build. In the official run, one negative tensor ascription in \nm{AmbientBoundary} produces a heartbeat timeout at the default limit. The same two inherited ascriptions were repeated with an enlarged limit, both producing type mismatches, not timeouts. Both supplementary author scripts also ran literally from the compiled independent tree. No exploratory failed compilation is counted as proof.
\paragraph{Reproducibility correction.}
The author's unchanged differential assembler correctly guards exact parent membership, but an ordinary checkout of the specified project contains 77 additional tracked top-level prose/freeze files even before auxiliary subdirectories. The native \nm{README.md} alone reproduces its rejection. To run the author script on a repository checkout, first project the merged subproject onto the exact 259 paths and verify every SHA-256 using \nm{project-parent.py} in the companion ZIP. The projected parent recovers all 295 target digests. Alternatively the review ZIP includes the complete verified source projection, offline manifest checker, and portable clean build/probe script. This is an invocation defect, not evidence of a different theorem.
\paragraph{Disposition.}
Accept the bounded pre-final four-to-two crossing and its consequences at the stated interface, with the explicit checkout projection in the reproduction instructions. This is one component of the manuscript pseudomonad construction, not closure of issue 105. The combination of a terminal-grade counterfeit that survives purity and an invertible crossed component that fails only terminal-payload purity identifies the cuts that a future simplification must still preserve.
\end{document}
